The combination of truncation-based hashing and exact-mode clusters enables handling very
large range queries that would be impractical for other approaches.

At $\mathcal{L} = 65536$, non-hybrid filters face a fundamental problem: the fingerprint
length $K = \lceil \log_2(n \cdot \mathcal{L} / \varepsilon) \rceil$ must grow by
$\log_2(\mathcal{L}) = 16$ bits to maintain the same FPR. For SODA, this means 16 extra
BPK. For Bloom, $\mathcal{L} = 65536$ point lookups per query is impractical.

Seg-ARE sidesteps this on two levels:

\paragraph{Exact clusters.} Each cluster stores its $[\min, \max]$ bounds. If a query
range $[a, b]$ overlaps a cluster's bounds, at least one stored key is inside $[a, b]$
--- that is a true positive, not a false positive. If $[a, b]$ does not overlap any
cluster, the cluster contributes nothing. If $[a, b]$ falls entirely within a cluster,
the sub-filter is exact (Section~\ref{sec:adaptive}), so false positives are impossible.
Either way, exact clusters produce zero false positives regardless of $\mathcal{L}$.

\paragraph{Truncation fallback.} The sparse fallback uses truncation hashing
(Section~\ref{sec:truncation}), where the phantom points of a stored key are confined
to the $2^t$-point prefix block containing that key
(Figure~\ref{fig:phantom-comparison}). A false positive can occur only when a query
\emph{endpoint} $a$ or $b$ falls inside such a block without containing the stored key
itself. Because these phantoms remain localized rather than scattering across the
universe as under SODA, the FPR of truncation is independent of the query range length
$\mathcal{L}$. The fallback thus avoids the $\log_2(\mathcal{L})$ BPK penalty that
SODA-based filters must pay.

\subsection{SOSD Results ($n = 2^{24}$, $\mathcal{L} = 65536$)}

\begin{figure}[h]
    \centering
    \includegraphics[width=0.9\textwidth]{fpr_bpk_books_L65536}
    \caption{FPR vs BPK --- SOSD Books (ISBN sequences, highly clustered),
        $n = 2^{24}$, $\mathcal{L} = 65536$. Seg-ARE achieves FPR~$< 10^{-8}$ at
        ${\sim}5$~BPK; all non-hybrid filters degrade to FPR~$\approx 1$.}
    \label{fig:books-l65536}
\end{figure}

\begin{figure}[h]
    \centering
    \includegraphics[width=0.9\textwidth]{fpr_bpk_wiki_L65536}
    \caption{FPR vs BPK --- SOSD Wiki (Wikipedia timestamps, moderately clustered),
        $n = 2^{24}$, $\mathcal{L} = 65536$. Seg-ARE dominates; only Grafite remains
        competitive at ${\sim}10$~BPK but with 2--3$\times$ more space.}
    \label{fig:wiki-l65536}
\end{figure}

On both datasets, Seg-ARE achieves FPR~$< 10^{-8}$ at ${\sim}5$~BPK, while Truncation,
SODA, and Bloom all degrade to FPR~$\approx 1$ at comparable budgets. SNARF plateaus at
FPR~$\sim 10^{-3}$. Only Grafite remains competitive on Wiki (${\sim}10$~BPK for
FPR~$10^{-2}$), but at 2--3$\times$ the space of Seg-ARE.

% TODO: add Facebook and OSM results when plots are regenerated
